\documentclass[11pt,a4paper]{article}
\usepackage{fontspec}
\usepackage{kotex}
\usepackage{geometry}
\usepackage{graphicx}
\usepackage{xcolor}
\usepackage{tcolorbox}
\usepackage{booktabs}
\usepackage{array}
\usepackage{longtable}
\usepackage{colortbl}
\usepackage{fancyhdr}
\usepackage{titlesec}
\usepackage{hyperref}
\usepackage{emoji}
\usepackage{amsmath}
\usepackage{amssymb}
\usepackage{enumitem}

\geometry{left=25mm, right=25mm, top=30mm, bottom=30mm}

\setmainfont{Noto Sans CJK KR}
\setsansfont{Noto Sans CJK KR}
\setmonofont{Noto Sans CJK KR}
\setemojifont{Noto Color Emoji}

\definecolor{primary}{HTML}{1976D2}
\definecolor{darkbrown}{HTML}{3E2723}
\definecolor{mediumbrown}{HTML}{5D4037}
\definecolor{lightbrown}{HTML}{8D6E63}
\definecolor{cream}{HTML}{FFF3E0}
\definecolor{lightgray}{HTML}{F5F5F5}
\definecolor{tableheader}{HTML}{E8E8E8}

\titleformat{\section}
  {\Large\bfseries\color{primary}}
  {\thesection.}{0.5em}{}
\titleformat{\subsection}
  {\large\bfseries\color{darkbrown}}
  {\thesubsection}{0.5em}{}
\titleformat{\subsubsection}
  {\normalsize\bfseries\color{mediumbrown}}
  {\thesubsubsection}{0.5em}{}

\renewcommand{\contentsname}{목차}
\renewcommand{\figurename}{그림}
\renewcommand{\tablename}{표}

\tcbset{
  tipbox/.style={colback=cream,colframe=primary,boxrule=1.5pt,arc=3mm,
    left=5mm,right=5mm,top=3mm,bottom=3mm,fonttitle=\bfseries\color{primary}},
  formulabox/.style={colback=lightgray,colframe=lightgray,boxrule=0pt,arc=2mm,
    left=5mm,right=5mm,top=3mm,bottom=3mm},
  summarybox/.style={colback=white,colframe=primary,boxrule=2pt,arc=4mm,
    left=5mm,right=5mm,top=5mm,bottom=5mm}
}

\hypersetup{colorlinks=true,linkcolor=primary,urlcolor=primary,bookmarks=true,unicode=true}


\pagestyle{fancy}
\fancyhf{}
\fancyhead[R]{\small\color{lightbrown}마이크로비트 Grove I2C 센서 시뮬레이터 이론해설서 v1.0.0.0}
\fancyfoot[C]{\thepage}
\renewcommand{\headrulewidth}{0.4pt}

\begin{document}

\begin{titlepage}
\centering
\vspace*{3cm}
{\Huge\bfseries\color{primary}\emoji{electric-plug} 마이크로비트 Grove I2C 센서 시뮬레이터}\\[0.5cm]
{\Huge\bfseries\color{primary} 이론해설서}
\\[1.5cm]
{\Large 시뮬레이터 버전: v9.2.0.0}
\\[2cm]
\begin{tcolorbox}[summarybox]
\centering
{\large 이 이론해설서는 마이크로비트 Grove I2C 센서 시뮬레이터이 다루는 핵심 개념과 원리를 설명합니다.}
\end{tcolorbox}
\vfill
{\large 사물인터넷과 빅데이터}\\[0.3cm]
{\large 청주교육대학교 컴퓨터교육과}
\vspace{1cm}
\end{titlepage}

\section*{1. 이 문서의 목적}

이 이론해설서는 \textbf{마이크로비트 Grove I2C 센서 시뮬레이터}이 다루는 핵심 개념과 원리를 설명합니다. 시뮬레이터를 사용하기 전에 배경 지식을 쌓거나, 사용 후 깊이 있는 이해를 위해 참조합니다.


\begin{tcolorbox}[summarybox]
\textbf{핵심 주제}: I2C 통신과 센서

\end{tcolorbox}


\section*{2. 핵심 개념}

\begin{center}
\renewcommand{\arraystretch}{1.4}
\begin{tabular}{|p{4.1cm}|p{10.9cm}|}
\hline
\rowcolor{tableheader}
\textbf{개념} & \textbf{설명} \\
\hline
\textbf{I2C 버스} & Inter-Integrated Circuit. 2개 선(SDA, SCL)으로 여러 장치를 연결하는 직렬 통신 프로토콜 \\
\hline
\textbf{I2C 주소} & 각 센서 고유 주소(0x00\textasciitilde{}0x7F). 하나의 버스에서 서로를 구분하는 '우편번호' \\
\hline
\textbf{Grove Shield} & 마이크로비트에 연결하여 I2C 포트를 6개로 확장하는 연결판 \\
\hline
\textbf{센서 융합} & 여러 센서 데이터를 결합하여 더 정확하거나 새로운 정보를 얻는 기법 \\
\hline
\textbf{디지털 센서 vs 아날로그 센서} & 디지털 센서는 보정된 수치를 직접 전송, 아날로그는 전압값 전송 후 변환 필요 \\
\hline
\end{tabular}
\end{center}


\section*{3. 원리 상세}

\subsection*{I2C 통신 원리}

I2C(Inter-Integrated Circuit)는 필립스(현 NXP)가 1982년에 개발한 직렬 통신 프로토콜입니다.


\textbf{핵심 특징:}

$\bullet$ \textbf{2선 구조}: SDA(데이터)와 SCL(클럭) 단 2개 선으로 통신

$\bullet$ \textbf{마스터-슬레이브}: 마이크로비트(마스터)가 센서(슬레이브)에게 명령

$\bullet$ \textbf{주소 기반}: 각 센서는 고유 주소를 가져 하나의 버스에 여러 센서 연결 가능

$\bullet$ \textbf{속도}: 표준 모드 100kHz, 고속 모드 400kHz


\textbf{통신 과정:}

1. 마스터가 START 신호 전송

2. 슬레이브 주소(7비트) + 읽기/쓰기 비트 전송

3. 해당 주소의 슬레이브가 ACK 응답

4. 데이터 전송 (8비트 단위)

5. 마스터가 STOP 신호로 통신 종료


\textbf{Grove 시스템의 장점:}

$\bullet$ 표준화된 4핀 커넥터 (VCC, GND, SDA, SCL)

$\bullet$ 꽂기만 하면 연결 완료 (납땜 불필요)

$\bullet$ 초등\textasciitilde{}대학생까지 쉽게 사용 가능



\section*{4. 실생활 응용 사례}

\textbf{사례 1}: 스마트 홈: 온습도·조도·CO\textsubscript{2} 센서로 실내 환경 자동 제어


\textbf{사례 2}: 스마트 농장: 토양 수분·온도·빛 센서로 작물 생장 모니터링


\textbf{사례 3}: 웨어러블: 가속도·심박·체온 센서 융합으로 건강 추적


\textbf{사례 4}: 기상관측소: 기압·온도·습도 데이터 수집 및 날씨 예보



\section*{5. 더 알아보기}

$\bullet$ \textbf{SPI vs I2C}: I2C보다 빠르지만 선이 4개 필요한 SPI 통신과의 비교

$\bullet$ \textbf{BLE 센서}: Bluetooth Low Energy로 무선 연결하는 방식

$\bullet$ \textbf{센서 캘리브레이션}: 센서 정확도를 높이기 위한 보정 과정

$\bullet$ \textbf{엣지 컴퓨팅}: 센서 데이터를 클라우드 대신 현장(엣지)에서 처리하는 방식



\section*{6. 핵심 용어 사전}

\begin{center}
\renewcommand{\arraystretch}{1.4}
\begin{tabular}{|p{3.4cm}|p{3.4cm}|p{8.2cm}|}
\hline
\rowcolor{tableheader}
\textbf{용어} & \textbf{학술 용어} & \textbf{쉬운 설명} \\
\hline
I2C 버스 & I2C 버스 & Inter-Integrated Circuit \\
\hline
I2C 주소 & I2C 주소 & 각 센서 고유 주소(0x00\textasciitilde{}0x7F) \\
\hline
Grove Shield & Grove Shield & 마이크로비트에 연결하여 I2C 포트를 6개로 확장하는 연결판 \\
\hline
센서 융합 & — & 여러 센서 데이터를 결합하여 더 정확하거나 새로운 정보를 얻는 기법 \\
\hline
디지털 센서 vs 아날로그 센서 & 디지털 센서 vs 아날로그 센서 & 디지털 센서는 보정된 수치를 직접 전송, 아날로그는 전압값 전송 후 변환 \\
\hline
\end{tabular}
\end{center}


\section*{7. 참고자료}

\begin{center}
\renewcommand{\arraystretch}{1.4}
\begin{tabular}{|p{1.5cm}|p{7.7cm}|p{5.8cm}|}
\hline
\rowcolor{tableheader}
\textbf{\#} & \textbf{자료명} & \textbf{비고} \\
\hline
1 & 마이크로비트Grove센서시뮬레이터\_퀵가이드 & 소프트웨어 사용 중 빠른 참조 \\
\hline
2 & 마이크로비트Grove센서시뮬레이터\_사용설명서 & 전체 기능 상세 \\
\hline
3 & 마이크로비트Grove센서시뮬레이터\_활용가이드 & 수업 적용 \\
\hline
4 & 마이크로비트Grove센서시뮬레이터\_개발참조 & 기술 사양 \\
\hline
\end{tabular}
\end{center}


\vfill
\begin{center}
\small\color{lightbrown}
\textit{마이크로비트 Grove I2C 센서 시뮬레이터 이론해설서}
\end{center}
\end{document}